\section{Correctness Properties}
\label{sec:rs-correctness}

This section establishes three correctness properties for the Recursive Spawning
mechanism: (1) drift prevention — the guarantee that spawned nodes cannot deviate
from the human principal's intent in undetectable or uncorrectable ways; (2) chain
integrity — the guarantee that the parent-pointer chain is tamper-evident and
reconstructable; and (3) termination — the guarantee that the escalation chain
from any node terminates at a human principal within a bounded number of steps.
Each property is derived from the three-layer integration specified in
Section~\ref{sec:rs-integration}.

% ---------------------------------------------------------------
\subsection{Drift Prevention}
\label{sec:rs-drift-prevention}

\textit{Drift} in a recursive spawning context refers to the phenomenon by which
a child node's behavior diverges from the principal's intent as encoded in the
parent's operating range — by expanding its own authority beyond what
$R(\text{parent})$ permits, by modifying the accountability chain, or by
replacing the human anchor with a machine actor as its escalation terminus.
Drift is the primary failure mode that the Recursive Spawning pillar is designed
to prevent, and its prevention depends on three structural properties acting in
concert.

\begin{theorem}[Drift Prevention]
\label{th:drift-prevention}
Let $T = (V, E)$ be a BX3 spawn tree satisfying all three constraints of
Section~\ref{sec:recursive-spawning-formal}: hierarchy finiteness, parent-pointer
immutability, and delegation scope inheritance.  Then for any node $v \in V$,
the accountability chain from $v$ to its human anchor $h(v)$ cannot be redirected,
truncated, or absorbed by any machine actor operating within the BX3 architecture.
\end{theorem}

\begin{proof}[Proof sketch]
The proof proceeds by establishing that three mutually reinforcing structural
properties jointly preclude each possible drift pathway.

\textit{Pathway 1: Parent-pointer modification.}  Suppose a machine actor $m$
attempts to redirect $\alpha(v)$ from its legitimate parent $p$ to a different
node $p'$.  By Invariant~\ref{inv:parent-pointer}, $\alpha(v)$ is a read-only
property of the \fact{} layer worksheet, enforced by the pre-commit hook.  Any
write attempt is intercepted and rejected before the state advances.  The rejection
event is recorded in the Forensic Ledger as
\texttt{parent-pointer-write-denied}$(v, p, p')$, and the Bailout Protocol is
triggered.  The drift attempt is both prevented and recorded.  No machine actor
has the authority to override this check.

\textit{Pathway 2: Human anchor replacement.}  Suppose a machine actor attempts
to replace $h(v)$ with a machine actor $m$ as the escalation terminus.  By
Invariant~\ref{inv:human-anchor}, $h(v)$ is established at node provisioning and
is never modified during the node's operational lifetime.  The human anchor is a
structural property of the \purpose{} layer, not a configurable parameter.  Any
attempt to modify it is rejected by the \fact{} layer's pre-commit hook, and the
Bailout Protocol is triggered.  The drift attempt is prevented.

\textit{Pathway 3: Authority expansion through spawning.}  Suppose a machine
actor at node $p$ attempts to spawn a child $c$ with $R(c) \not\subset R(p)$,
thereby expanding delegated authority beyond what the principal sanctioned.  By
Invariant~\ref{inv:scope-inheritance}, the \bounds{} layer evaluates
$R(c) \subset R(p)$ before every spawn event.  If the check fails, the
$e_{\text{scope-violation}}$ event is emitted, the spawning state machine
transitions to \texttt{SPAWN\_LOCKED}, and the Bailout Protocol is triggered.
The attempted authority expansion is prevented and recorded.

\textit{Pathway 4: Chain truncation via non-termination.}  Suppose a machine
actor attempts to prevent the escalation chain from reaching $h(v)$ by causing
the chain to loop or diverge.  By hierarchy finiteness (Invariant
~\ref{inv:hierarchy-finiteness}), the parent-pointer chain from any node has
maximum length $D_{\max}$, and by Invariant~\ref{inv:spawn-accountability} it
terminates at a human principal.  The chain cannot loop because $\alpha(v)$ is
immutable and acyclic (a node cannot be its own ancestor).  The chain cannot
terminate at a machine actor because Invariant~\ref{inv:human-anchor} forbids
the human anchor from being replaced.  The drift attempt is structurally
impossible.

Since all four drift pathways are precluded, the accountability chain from any
node $v$ to its human anchor $h(v)$ is non-divertable.  $\square$
\end{proof}

\medskip
\noindent\textbf{Role of the Forensic Ledger.}  Even though pre-commit hooks
prevent drift at runtime, the Forensic Ledger ensures that any attempted drift —
successful or not — is permanently recorded with cryptographic chain integrity.
Drift attempts are therefore not merely prevented but also attributable: the
identity of the node that attempted the drift, the time of the attempt, and the
nature of the violation are all reconstructable by any auditor.  The combination
of prevention and attribution is what makes the upstream accountability guarantee
operational in high-stakes deployments.

The drift prevention theorem depends critically on the \textit{conjunction} of
all three constraints.  No single constraint is sufficient: parent-pointer
immutability alone does not prevent authority expansion; delegation scope
inheritance alone does not prevent chain redirection; hierarchy finiteness alone
does not prevent the human anchor from being replaced.  It is the conjunction of
all three, enforced by all three layers operating in concert, that yields the
non-divertability guarantee.

% ---------------------------------------------------------------
\subsection{Chain Integrity}
\label{sec:rs-chain-integrity}

\textit{Chain integrity} refers to the guarantee that the parent-pointer chain
from any node to its human anchor is complete, tamper-evident, and reconstructable.
The chain is the escalation path used by the Bailout Protocol; if the chain is
incomplete or tampered with, the escalation path is unreliable and the Liveness
guarantee is compromised.

\begin{property}[Chain Integrity]
\label{prop:chain-integrity}
For any node $v$ in a BX3 spawn tree, the accountability chain
$C(v) = \langle v, \alpha(v), \alpha^{(2)}(v), \ldots, \alpha^{(k)}(v) \rangle$
where $\alpha^{(k)}(v) = h(v)$ satisfies:
\begin{enumerate}
\item $k \leq D_{\max}$ (finite length);
\item $\alpha^{(i)}(v)$ is a machine node for all $0 \leq i < k$;
\item $\alpha^{(k)}(v)$ is a human principal;
\item For all $0 \leq i < j \leq k$, $\alpha^{(i)}(v) \neq \alpha^{(j)}(v)$ (acyclic);
\item Every entry in $C(v)$ appears as a \texttt{child-provisioned} event in the
Forensic Ledger with matching $\alpha(\cdot)$ and $h(\cdot)$ values.
\end{enumerate}
\end{property}

Property~\ref{prop:chain-integrity} is established by construction and verified by
the \fact{} layer at each spawn event.  Items (1)–(3) follow from the hierarchy
finiteness constraint, the delegation scope inheritance constraint, and the human
anchor immutability invariant.  Item (4) follows from parent-pointer immutability:
since $\alpha(v)$ is established once at node creation and never modified, and
since $\alpha(v) \neq v$ for all $v$, the parent-pointer relation defines a
strict partial order with no cycles.  Item (5) follows from the Forensic Ledger
completeness invariant: every spawn event is recorded before subsequent actions
are permitted, and every \texttt{child-provisioned} entry encodes both
$\alpha(c) = p$ and $h(c) = h(p)$.

The tamper-evidence property follows from the Forensic Ledger's SHA-256 hash
chaining.  Each entry carries the hash of the preceding entry.  Any modification
to an entry — node identifier, event type, context, or parent pointer — breaks
the chain at the point of modification.  The break is detectable by recomputing
the chain of hashes from the genesis entry to the entry in question and comparing
the recomputed hash against the stored hash.  An auditor without access to the
ledger's contents can detect tampering with high confidence.

Reconstructability follows from items (4)–(5) and the cryptographic chaining.
Given any node $v$, an auditor reconstructs $C(v)$ by querying the Forensic
Ledger for all \texttt{child-provisioned} entries involving $v$, extracting the
$\alpha(v)$ field from each, and traversing the chain until reaching a node that
does not appear as a child in any \texttt{child-provisioned} entry — that node
is $h(v)$.  The reconstruction is deterministic and reproducible; any two
auditors who independently reconstruct $C(v)$ obtain identical results.

% ---------------------------------------------------------------
\subsection{Termination}
\label{sec:rs-termination}

\textit{Termination} is the guarantee that the escalation chain from any node
always reaches a human principal in finite time, and that the Bailout Protocol's
escalation algorithm terminates.  This is the structural foundation of the BX3
Framework's Liveness guarantee.

\begin{property}[Finite Termination]
\label{prop:finite-termination}
For any node $v$ in a BX3 spawn tree, the escalation chain $C(v)$ terminates
at a human principal $h(v)$ in at most $D_{\max}$ steps.
\end{property}

\begin{proof}[Proof sketch]
By the hierarchy finiteness constraint, every node $v$ has depth $d(v) \leq D_{\max}$,
where $d(\text{root}) = 0$ and $d(v) = d(\alpha(v)) + 1$ for all $v \neq \text{root}$.
The root node is provisioned directly by the human principal and therefore satisfies
$h(\text{root}) = \text{principal}(\text{root})$.  For any node $v$, the chain
$v, \alpha(v), \alpha^{(2)}(v), \ldots, \alpha^{(d(v))}(v)$ exists and has length
$d(v) \leq D_{\max}$.  By Invariant~\ref{inv:human-anchor}, $\alpha^{(d(v))}(v)$
is a human principal.  Therefore $C(v)$ terminates at a human principal in at most
$D_{\max}$ steps.  $\square$
\end{proof}

Property~\ref{prop:finite-termination} is critical for the Bailout Protocol's
operational correctness.  The Bailout Protocol traverses $C(v)$ from $v$ upward,
escalating exceptions at each node until reaching $h(v)$.  If $C(v)$ were infinite
or non-terminating, the escalation would stall — the Bailout Protocol would never
reach the human principal, and the Liveness guarantee would be vacuous.  The
property ensures that the escalation algorithm always completes in at most
$D_{\max}$ steps.

\begin{property}[Deterministic Escalation Path]
\label{prop:deterministic-escalation}
For any node $v$ and any exception state $e$ at $v$, the Bailout Protocol's
escalation path from $v$ to $h(v)$ is uniquely determined and is independent of
the content of $e$ or the state of the \bounds{} layer at any node on the path.
\end{property}

Property~\ref{prop:deterministic-escalation} follows from the immutability of
the parent-pointer chain.  The escalation algorithm traverses
$\alpha(v), \alpha^{(2)}(v), \ldots, \alpha^{(k)}(v)$ where $k \leq D_{\max}$.
Each $\alpha^{(i)}(v)$ is immutable and deterministic — established once at node
creation and never modified.  The path does not depend on the exception type, the
exception severity, or any runtime state.  It is a structural property of the
spawn tree, not an operational decision made by any machine actor.  This means
that the escalation path is predictable and auditable before any exception occurs:
the human anchor can determine the complete escalation path for any node by
reading the Forensic Ledger entries for all \texttt{child-provisioned} events
along the chain.

% ---------------------------------------------------------------
\subsection{Collective Correctness}
\label{sec:rs-collective-correctness}

Together, drift prevention, chain integrity, and finite termination constitute the
full correctness specification for Recursive Spawning.  They establish that a BX3
spawn tree is a structurally sound representation of the delegation chain: every
edge represents a genuine, auditable, bounded delegation of authority from a
principal (human or machine) to a child node, and the chain from any leaf node to
its human anchor is fully recorded, finite in depth, tamper-evident, and
non-circumventable.

These correctness properties are not conditional on the honesty or correctness of
any machine actor.  They are enforced by the structural properties of the
three-layer architecture — specifically by the pre-commit hooks of the \fact{}
layer — and they hold regardless of what any agent or node attempts to do at
runtime.  This distinguishes the BX3 Framework's correctness guarantees from
conventional audit-based accountability systems: the guarantees are architectural,
not operational.
